%%%% RESUMO
%%
%% Apresentação concisa dos pontos relevantes de um texto, fornecendo uma visão rápida e clara do conteúdo e das conclusões do
%% trabalho.

\begin{resumoutfpr}%% Ambiente resumoutfpr
A seleção automática de pipelines de aprendizado de máquina é um problema central em AutoML, pois o desempenho dos modelos depende das etapas de pré-processamento e da combinação de algoritmos empregados. Este trabalho propõe uma abordagem de meta-aprendizado para recomendação e ranqueamento de pipelines de classificação baseada em uma representação estatística distribucional de conjuntos de dados, sem treinamento neural. Na Fase 1, construiu-se um meta-conjunto com 201 conjuntos obtidos do OpenML e 60 pipelines formados por combinações de escalonadores, seletores de atributos e classificadores, avaliados por validação cruzada estratificada. Cada conjunto foi representado por um vetor estatístico de 64 dimensões, inspirado nos critérios de representação do Dataset2Vec. Na Fase 2, um preditor k-NN multi-saída com estrutura KDTree foi avaliado por Leave-One-Out (LOO), utilizando Spearman e Kendall Tau para correlação global do ranking, e Kendall Tau ponderado e NDCG como métricas complementares voltadas ao topo da recomendação. Para $k=10$, o preditor obteve Spearman médio de $0{,}58 \pm 0{,}36$, Kendall de $0{,}48 \pm 0{,}34$ e Kendall ponderado de $0{,}52 \pm 0{,}35$, indicando capacidade de capturar a ordem relativa dos pipelines, sobretudo nas primeiras posições. Uma demonstração com conjunto de dados retido é apresentada apenas como exemplo operacional do recomendador, sem substituir a validação estatística por LOO. A correlação entre dimensões da representação e Kendall ponderado ($p < 0{,}0001$) sugere que o vetor estatístico captura estrutura relevante para a recomendação. Os resultados indicam que representações estatísticas aliadas a k-NN constituem uma abordagem viável e reprodutível para recomendação de pipelines em AutoML.
\end{resumoutfpr}
